% Generated draft: requires mathematical and source review.
\documentclass[11pt]{article}
\usepackage[margin=25mm]{geometry}
\usepackage{fontspec}
\setmainfont{Latin Modern Roman}
\usepackage{amsmath,amssymb,mathtools,unicode-math}
\setmathfont{Latin Modern Math}
\usepackage{enumitem,needspace,xurl,hyperref}
\setlength{\parindent}{0pt}
\setlength{\parskip}{5pt}
\setlength{\emergencystretch}{3em}
\newcommand{\N}{\mathbb{N}}
\newcommand{\Z}{\mathbb{Z}}
\newcommand{\Q}{\mathbb{Q}}
\newcommand{\R}{\mathbb{R}}
\newcommand{\C}{\mathbb{C}}
\newcommand{\F}{\mathbb{F}}
\newcommand{\A}{\mathbb{A}}
\newcommand{\PP}{\mathbb{P}}
\newcommand{\E}{\mathbb{E}}
\newcommand{\eps}{\varepsilon}
\newcommand{\dd}{\,\mathrm{d}}
\begin{document}
\section*{UnsolvedMath ID: 30005026}
\textbf{Source problem code:} catalogue entry
\subsection{Definitions and mathematical statement}
The supplied question concerns simultaneous spatial and informational efficiency of a finitary coding. Its ambient index set, translation action, metric, state spaces, and entropy convention are not specified. The following is a formulation template relative to those choices, rather than an assertion that any particular choices were supplied.

Fix an ambient process setting with a distinguished site \(o\), finite metric balls \(B(o,r)\), and translations. A factor map \(\Phi\) is a measurable map commuting with these translations; saying that \(X\) is a factor of \(Y\) means \(\Phi(Y)\stackrel{\mathrm d}=X\). The coding is finitary if, almost surely, the output \(\Phi(Y)_o\) is determined by the input on \(B(o,R)\) for some finite random coding radius \(R\). Exponentially decaying coding-radius tails mean that there are constants \(C,c>0\) such that
\[
\Pr(R>r)\le C e^{-cr}\qquad\text{for every integer }r\ge0.
\]
A finite-valued input process has a finite alphabet. Write \(h(Z)\) for the intended entropy per site of a process \(Z\), once its definition and normalization have been fixed.

\textbf{Question, subject to the unresolved qualifications below.} Suppose there exist an independent and identically distributed input process \(Y^{(0)}\) and a finitary factor map \(\Phi_0\) such that
\[
\Phi_0(Y^{(0)})\stackrel{\mathrm d}=X,
\qquad
\Pr(R_0>r)\le C_0e^{-c_0r}
\quad(r\in\Z_{\ge0})
\]
for some \(C_0,c_0>0\). Must it be true that, for every \(\varepsilon>0\), there exist a finite alphabet \(A_\varepsilon\), an independent input process \(Y^{(\varepsilon)}\) taking values in \(A_\varepsilon\), and a finitary factor map \(\Phi_\varepsilon\) such that
\[
\Phi_\varepsilon(Y^{(\varepsilon)})\stackrel{\mathrm d}=X,
\qquad
\Pr(R_\varepsilon>r)\le C_\varepsilon e^{-c_\varepsilon r}
\quad(r\in\Z_{\ge0})
\]
for some \(C_\varepsilon,c_\varepsilon>0\), while the input entropy is arbitrarily close to that of \(X\)? If the intended entropies are finite real entropy rates, the closeness requirement can be written
\[
\bigl|h(Y^{(\varepsilon)})-h(X)\bigr|<\varepsilon.
\]
In a setting where factor maps cannot increase this entropy, this is equivalently
\[
h(X)\le h(Y^{(\varepsilon)})<h(X)+\varepsilon.
\]
The constants and coding may depend on \(\varepsilon\); the supplied statement imposes no uniform exponential bound across \(\varepsilon\).

The record does not specify whether the second input must also be identically distributed, whether \(h(X)<\infty\) is assumed, or which ambient setting makes the entropy comparison meaningful. These omissions prevent an unconditional precise formulation. The verbatim extracted question refers only to ``both properties''; their identification with the two requirements above comes from the supplied reconstructed statement.
\subsection{Short English statement}
Given a process obtained from an i.i.d. input by a finitary coding with exponentially decaying radius tails, can one simultaneously retain exponential tails and use a finite-valued independent input whose entropy is arbitrarily close to the output entropy? The ambient process setting, entropy assumptions, and whether the new input must be i.i.d. remain unspecified. The supplied record labels the problem open; this draft does not independently verify that status.
\subsection{Sources}
\begin{itemize}
\item Local registry: \texttt{data/unsolvedmath/problems.json}, with \texttt{catalog/problems.json} as the project supplementary registry; supplied permanent UnsolvedMath ID \(30005026\), title ``Simultaneous Spatial and Informational Coding Efficiency,'' source problem code \texttt{catalogue entry}. Registry files were not independently inspected for this draft.
\item Supplied report citation: \emph{Mini-Workshop: Descriptive Combinatorics, LOCAL Algorithms and Random Processes} (2023), \emph{Oberwolfach Reports}, pp.\ 429--455, DOI \texttt{10.4171/owr/2022/8}. Supplied source URL: \url{https://doi.org/10.4171/owr/2022/8}. The original extraction includes the page marker 449 and asks, ``Is it the case that both properties can be realized simultaneously?''
\item Supplied literature source, labeled ``OWR finitary-process report'': \url{https://doi.org/10.4171/owr/2022/8}.
\item Literature assessment dated exactly \texttt{2026-08-22}: the supplied curation labels the question open and states that the report gives examples and separate partial results while leaving simultaneous compatibility open. This is supplied dated curation, not independently verified literature assessment; time-sensitive status and the cited claims require independent verification.
\item The record labels its reconstructed statement \texttt{corrected\_verified} and reports that it was checked against the official Oberwolfach report. That verification claim is supplied metadata, not verification performed here. The record gives proposed year 2022 and an \texttt{owr-year:2022} tag, whereas its report citation gives 2023. No individual proposer is supplied; attribution is to the report and its workshop open-problem discussion.
\end{itemize}
\end{document}
